% discussion.tex -- Discussion and Conclusion (merged to save space)

The experiments validated the core claim: Langevin dynamics on the modern Hopfield energy converts deterministic attention into a principled stochastic sampler governed by a single temperature parameter. The synthetic experiments confirmed a smooth phase transition between retrieval and generation (Fig.~\ref{fig:experiments}a), convergence to the correct Boltzmann target (Fig.~\ref{fig:experiments}b), and graceful degradation with load ratio (Fig.~\ref{fig:phase-diagram}). The strongest result came from protein sequences (Table~\ref{tab:protein-main}), where the training-free score function preserved family-level fidelity at position-specific and co-evolutionary levels that the VAE could not maintain from $K{=}68$ examples. On MNIST, SA led the best learned baseline by $2.6{\times}$ on novelty and $2.0{\times}$ on diversity, while DDPM failed across all memory sizes tested. The entropy inflection (Proposition~\ref{prop:entropy-derivative}) located the retrieval-to-generation boundary from data alone, correctly predicting $\beta^*{=}3.85$ on RRM, and the masking extension lifted SA to $96.0\%$ subject-conditional accuracy on Olivetti from a $10.4\%$ unconditional baseline.

A deliberate choice shapes our experimental design: we evaluate at memory sizes $K$ from 68 to 3{,}500, the low-to-moderate data regime where learned generative models are most constrained, matching most real low-data domains (small protein families, rare-event catalogs, specialized image collections, small clinical cohorts). Score-based and diffusion models require enough data to learn an accurate score function, ill-posed at small $K$ as DDPM's failure illustrates; stochastic attention sidesteps this bottleneck because $\nabla\log p_\beta = \beta(\mathbf{T}-\boldsymbol{\xi})$ is exact for any $K$, reusing the primitives of a standard attention head with no learned score network, training loop, or contrastive objective. The Hopfield energy yields convergence guarantees in the convex regime (Corollary~\ref{cor:convergence}) and geometric ergodicity for all $\beta>0$, while at high $\beta$ the multi-chain protocol works in practice. The same closed-form score function generalizes to soft-bias conditioning ($b_{k}\in\mathbb{R}$), recovering the multiplicity-weight conditioning of Varner~\cite{varnerConditioningProtein2026}; concurrent work extends the framework to eight Pfam families~\cite{varnerTrainingFreeProtein2026} and validated synthetic patient generation from small clinical cohorts~\cite{varnerSyntheticPatient2026}, and integration with pretrained multi-head attention is future work (Appendix~\ref{app:multihead}).

\paragraph{Limitations.}
ULA discretization bias scales with $\alpha$; MALA (Appendix~\ref{app:mala-algorithm}) is preferred at larger step sizes, though the $99.2\%$ acceptance rate at $\alpha{=}0.01$ confirms negligible bias here. At high temperature ($\beta{=}200$) SA samples are structured interpolations rather than crisp novel images. Inter-basin mixing at high $\beta$ grows exponentially with barrier height; quantitative scaling laws remain open, and the multi-chain protocol is presently the only practical workaround. The method requires $\mathbf{X}$ to be known and fixed; streaming-memory or latent-memory settings would require extension, as would composing SA with learned multi-head attention in full transformers.
